\section{Related literature and model provenance}
\label{sec:literature}

The intellectual ancestry runs from the two-period consumption-loan and neoclassical OLG economies of \citet{Samuelson1958} and \citet{Diamond1965} to the 55-period computational life-cycle model of \citet{AuerbachKotlikoff1987}, which established the OLG model as the workhorse for quantitative fiscal-policy analysis --- tax reform, social security, and generational incidence. Modern large-scale descendants add idiosyncratic ability heterogeneity, realistic demographics, and detailed fiscal institutions; \citet{NishiyamaSmetters2007} is a prominent example in the social-security literature, and the OBR has recently built its own UK OLG model in the same tradition \citep{OBR2025OLG}.

OG-Core is the open-source embodiment of this tradition maintained by the Policy Simulation Library community. Its theory, calibration interface, solution methods, and Python API are documented in \citet{DeBackerEvans2024}, which serves as the model's canonical reference; \citet{EvansDeBacker2024} provide a textbook treatment of building and solving OLG models for policy analysis that covers most of the structures in OG-Core, and \citet{EvansPhillips2014} analyse the time-path-iteration solution method the framework uses for transition dynamics. The framework separates \emph{theory} (OG-Core, a country-agnostic system of equations and solvers) from \emph{calibration} (country packages that supply demographics, tax detail, and macro parameters). OG-USA, the United States calibration, is documented in \citet{DeBackerEvansOGUSA} and has been used for dynamic scoring of US federal tax reform; OG-UK \citep{OGUKdocs} is the United Kingdom calibration, originally built on OpenFisca-UK and now on PolicyEngine UK \citep{PolicyEngineUK}, sourcing calibration targets from the ONS, the OBR, the Bank of England, and GOV.UK.

The methodological bridge that makes the model usable for statutory reform scoring is \citet{DeBackerEvansPhillips2019}, who propose estimating parametric effective- and marginal-tax-rate functions --- jointly in labour and capital income --- from the output of a tax microsimulation model, and embedding those functions in the dynamic general-equilibrium model's household problem. The functional-form menu draws on the empirical tax-function literature: the three-parameter total-income schedule of \citet{GouveiaStrauss1994} and the log-linear progressivity form of \citet{HSV2017}, alongside DeBacker--Evans--Phillips's own richer bivariate specification. PolicyEngine Macro's contribution is not to this theory but to the pipeline: replacing the microsimulation input with PolicyEngine's open-source UK model over enhanced FRS microdata, and exposing the whole loop --- reform in, steady-state general-equilibrium impact out --- behind the suite's common scoring interface.
